\documentclass[12pt,letterpaper]{article}
\usepackage{graphicx,textcomp}
\usepackage{natbib}
\usepackage{setspace}
\usepackage{fullpage}
\usepackage{color}
\usepackage[reqno]{amsmath}
\usepackage{amsthm}
\usepackage{fancyvrb}
\usepackage{amssymb,enumerate}
\usepackage[all]{xy}
\usepackage{endnotes}
\usepackage{lscape}
\newtheorem{com}{Comment}
\usepackage{float}
\usepackage{hyperref}
\newtheorem{lem} {Lemma}
\newtheorem{prop}{Proposition}
\newtheorem{thm}{Theorem}
\newtheorem{defn}{Definition}
\newtheorem{cor}{Corollary}
\newtheorem{obs}{Observation}
\usepackage[compact]{titlesec}
\usepackage{dcolumn}
\usepackage{tikz}
\usetikzlibrary{arrows}
\usepackage{multirow}
\usepackage{xcolor}
\newcolumntype{.}{D{.}{.}{-1}}
\newcolumntype{d}[1]{D{.}{.}{#1}}
\definecolor{light-gray}{gray}{0.65}
\usepackage{url}
\usepackage{listings}
\usepackage{color}

\definecolor{codegreen}{rgb}{0,0.6,0}
\definecolor{codegray}{rgb}{0.5,0.5,0.5}
\definecolor{codepurple}{rgb}{0.58,0,0.82}
\definecolor{backcolour}{rgb}{0.95,0.95,0.92}

\lstdefinestyle{mystyle}{
	backgroundcolor=\color{backcolour},   
	commentstyle=\color{codegreen},
	keywordstyle=\color{magenta},
	numberstyle=\tiny\color{codegray},
	stringstyle=\color{codepurple},
	basicstyle=\footnotesize,
	breakatwhitespace=false,         
	breaklines=true,                 
	captionpos=b,                    
	keepspaces=true,                 
	numbers=left,                    
	numbersep=5pt,                  
	showspaces=false,                
	showstringspaces=false,
	showtabs=false,                  
	tabsize=2
}
\lstset{style=mystyle}
\newcommand{\Sref}[1]{Section~\ref{#1}}
\newtheorem{hyp}{Hypothesis}


\title{Problem Set 4}
\date{Due: November 26, 2026}
\author{Owen Stidman}


\begin{document}
	\maketitle
	\section*{Instructions}
	\begin{itemize}
		\item Please show your work! You may lose points by simply writing in the answer. If the problem requires you to execute commands in \texttt{R}, please include the code you used to get your answers. Please also include the \texttt{.R} file that contains your code. If you are not sure if work needs to be shown for a particular problem, please ask.
		\item Your homework should be submitted electronically on GitHub.
		\item This problem set is due before 23:59 on Thursday November 26, 2026. No late assignments will be accepted.
	\end{itemize}



	\vspace{.5cm}
\section*{Question 1: Economics}
\vspace{.25cm}
\noindent 	
In this question, use the \texttt{prestige} dataset in the \texttt{car} library. First, run the following commands:

\begin{verbatim}
install.packages(car)
library(car)
data(Prestige)
help(Prestige)
\end{verbatim} 


\noindent We would like to study whether individuals with higher levels of income have more prestigious jobs. Moreover, we would like to study whether professionals have more prestigious jobs than blue and white collar workers.

\newpage
\begin{enumerate}
	
	\item [(a)]
	Create a new variable \texttt{professional} by recoding the variable \texttt{type} so that professionals are coded as $1$, and blue and white collar workers are coded as $0$ (Hint: \texttt{ifelse}). 	\vspace{.5cm}
	
	Using \textit{ifelse()} and the code below, I first create a copy of the dataset to avoid overwriting the original data, then create a new numeric variable called \textit{professional} where the value "prof" (professional) from the categorical \textit{type} variable is coded as 1 and the values "wc" and "bc" (white- and blue-collar) are both coded as zero. The discrete numeric \textit{professional} variable is a dummy variable containing 102 observations. 
	
	\lstinputlisting[language=R, firstline=46, lastline=48]{PS04_answers_OS.R} 
	
	\vspace{6cm}
	
	
	\item [(b)]
	Run a linear model with \texttt{prestige} as an outcome and \texttt{income}, \texttt{professional}, and the interaction of the two as predictors (Note: this is a continuous $\times$ dummy interaction.) \vspace{.5cm}
	
	Using the below code, I run a linear regression with \texttt{prestige} as an outcome and \texttt{income}, \texttt{professional}, and the interaction of the two as predictors.
	
		\lstinputlisting[language=R, firstline=54, lastline=64]{PS04_answers_OS.R} 
	
	The below table (Table 1) contains a summary of the model.
		
		\begin{table}[!htbp] \centering 
			\caption{} 
			\label{} 
			\begin{tabular}{@{\extracolsep{5pt}}lc} 
				\\[-1.8ex]\hline 
				\hline \\[-1.8ex] 
				& \multicolumn{1}{c}{\textit{Dependent variable:}} \\ 
				\cline{2-2} 
				\\[-1.8ex] & prestige \\ 
				\hline \\[-1.8ex] 
				income & 0.003$^{***}$ \\ 
				& (0.0005) \\ 
				& \\ 
				professional & 38.120$^{***}$ \\ 
				& (4.080) \\ 
				& \\ 
				professional & $-$0.002$^{***}$ \\ 
				& (0.001) \\ 
				& \\ 
				Constant & 20.804$^{***}$ \\ 
				& (2.539) \\ 
				& \\ 
				\hline \\[-1.8ex] 
				Observations & 102 \\ 
				R$^{2}$ & 0.789 \\ 
				Adjusted R$^{2}$ & 0.783 \\ 
				Residual Std. Error & 8.018 (df = 98) \\ 
				F Statistic & 122.338$^{***}$ (df = 3; 98) \\ 
				\hline 
				\hline \\[-1.8ex] 
				\textit{Note:}  & \multicolumn{1}{r}{$^{*}$p$<$0.1; $^{**}$p$<$0.05; $^{***}$p$<$0.01} \\ 
			\end{tabular} 
		\end{table} 
	
	The table is also included in the R output below:
	
	\begin{verbatim}
		Call:
		lm(formula = prestigecopy$prestige ~ prestigecopy$income + prestigecopy$professional + 
		(prestigecopy$income * prestigecopy$professional), data = prestigecopy)
		
		Residuals:
		Min      1Q  Median      3Q     Max 
		-14.854  -5.429  -1.055   4.513  29.762 
		
		Coefficients:
		Estimate Std. Error t value Pr(>|t|)    
		(Intercept)                                   20.8035385  2.5387225   8.194 9.74e-13 ***
		prestigecopy$income                            0.0032691  0.0004565   7.161 1.49e-10 ***
		prestigecopy$professional                     38.1200003  4.0797620   9.344 3.22e-15 ***
		prestigecopy$income:prestigecopy$professional -0.0024239  0.0005304  -4.570 1.42e-05 ***
		---
		Signif. codes:  0 ‘***’ 0.001 ‘**’ 0.01 ‘*’ 0.05 ‘.’ 0.1 ‘ ’ 1
		
		Residual standard error: 8.018 on 98 degrees of freedom
		Multiple R-squared:  0.7893,	Adjusted R-squared:  0.7828 
		F-statistic: 122.3 on 3 and 98 DF,  p-value: < 2.2e-16  
	\end{verbatim}
	
	The intercept coefficient is 20.8035385. On average, holding \textit{income} and \textit{professional} constant at zero (i.e., a non-professional worker earning no income) would be expected to have a "prestige" score of 20.8035385. 
	
	The slope coefficient for \textit{income} is 0.0032691. On average, holding professional job status constant, a one-Canadian-dollar increase in income would be associated with a predicted prestige score increase of 0.0032691. 
	
	The slope coefficient for \textit{professional} is 38.1200003. Non-professional (i.e., white- and blue-collar) is the reference category. In comparison with non-professional jobs, holding income constant, professionals have, on average, a prestige score that is higher by 38.1200003. 
	
	The slope coefficient for \textit{(income x professional)} (i.e., the interation between the \textit{income} and \textit{professional}) variables, is -0.0024239. In comparison with non-professional jobs, on average, the effect of a one-Canadian-dollar increase in income on prestige is -0.0024239 units lower for professionals. 
	
	The adjusted R-squared value is 0.7828. Approximately 78.28 percent of the variation in prestige is explained by the variation in income, professional status, and the interaction between income and professional status. 
	
	\vspace{6cm}
	\item [(c)]
	Write the prediction equation based on the result. \vspace{.5cm}
	
	Table 1 and the summary output above contain the coefficients for the regression equation below. \vspace{.5cm}
	
\textbf{General model:}
	
	$\hat{prestige}$ = 20.8035385 + 0.0032691\textit{income} + 38.1200003\textit{professional} - 0.0024239\textit{(income x professional)} + $\epsilon$
	
	\vspace{.5cm} As the value of the \textit{professional} variable changes both the intercept and the slope coefficient for \textit{income}, there are different models for professionals and for non-professionals. 	\vspace{.5cm}
	
\textbf{Model for professionals (\textit{professional = 1}):}
	
	$\hat{prestige}$ = 20.8035385 + 0.0032691\textit{income} + (38.1200003*1) - 0.0024239\textit{(income * 1)} + $\epsilon$
	
	Using basic arithmetic, this simplifies to:
	
	$\hat{prestige}$ = 58.9235388 + 0.0008452\textit{income} + $\epsilon$
	
	On average, for professionals, a one-Canadian-dollar increase in income is associated with a predicted increase in prestige of 0.0008452.
	
		\vspace{.5cm} \textbf{Model for non-professionals (\textit{professional = 0}):} 
	
	$\hat{prestige}$ = 20.8035385 + 0.0032691\textit{income} + (38.1200003*0) - 0.0024239\textit{(income * 0)} + $\epsilon$
	
	Using basic arithmetic, this simplifies to:
	
	$\hat{prestige}$ = 20.8035385 + 0.0032691\textit{income} + $\epsilon$
	
		On average, for non-professionals, a one-Canadian-dollar increase in income is associated with a predicted increase in prestige of 0.0032691.

	
\newpage
	\item [(d)]
	Interpret the coefficient for \texttt{income}. \vspace{.5cm}
	
	Using the general model above:
	
	On average, holding professional job status constant, a one-Canadian-dollar increase in income would be associated with a predicted prestige score increase of 0.0032691.  \vspace{.5cm}
	
	Using the model for professionals (\textit{professional = 1}) above:
	
	On average, for professionals, a one-Canadian-dollar increase in income is associated with a predicted increase in prestige of 0.0008452. \vspace{.5cm}
	
	Using the model for non-professionals (\textit{professional = 0}) above:
	
	On average, for non-professionals, a one-Canadian-dollar increase in income is associated with a predicted increase in prestige of 0.0032691.
	
	\vspace{5cm}	
	\item [(e)]
	Interpret the coefficient for \texttt{professional}. \vspace{.5cm}
	
	As \textit{professional} takes a value of 1 or 0 for the professional and non-professional models, the interpretation of this coefficient only applies to the general model. 
	
	Non-professional (i.e., white- and blue-collar) is the reference category. In comparison with non-professional jobs, holding income constant, professionals have, on average, a prestige score that is higher by 38.1200003. 
	
	\newpage
	\item [(f)]
	What is the effect of a \$1,000 increase in income on prestige score for professional occupations? In other words, we are interested in the marginal effect of income when the variable \texttt{professional} takes the value of $1$. Calculate the change in $\hat{y}$ associated with a \$1,000 increase in income based on your answer for (c). \vspace{.5cm}	
	
	The expected prestige score for a professional earning no income (i.e., \textit{income = 0} and \textit{professional} = 1) is:
	
	$\hat{prestige}$ = 20.8035385 + (0.0032691*0) + (38.1200003*1) - 0.0024239(0*1) + $\epsilon$
	= 20.8035385 + 0 + 38.1200003 - 0
	= 20.8035385 + 38.1200003 
	= 58.9235388.  \vspace{.5cm}	
	
	The expected prestige score for a professional earning 1000 Canadian dollars (i.e., \textit{income = 1000} and \textit{professional} = 1) is:
	
	$\hat{prestige}$ = 20.8035385 + (0.0032691*1000) + (38.1200003*1) - 0.0024239(1000*1) + $\epsilon$
	= 20.8035385 + (0.0032691*1000) + (38.1200003*1) - 0.0024239(1000*1) 
	= 20.8035385 + 3.2691 + 38.1200003 - 2.4239
	= 59.7687388. \vspace{.5cm}	
	
	59.7687388 - 58.9235388 = \textbf{0.8452}
	
	On average, a 1000-Canadian-dollar increase in income would be associated with an expected prestige score increase of 0.8452.
	
	\vspace{11cm}
	
	
	\item [(g)]
	What is the effect of changing one's occupations from non-professional to professional when her income is \$6,000? We are interested in the marginal effect of professional jobs when the variable \texttt{income} takes the value of $6,000$. Calculate the change in $\hat{y}$ based on your answer for (c). \vspace{1cm}
	
	The expected prestige score for a non-professional earning 6000 Canadian dollars (i.e., \textit{professional=0} and \textit{income = 6000}) is:
	
		$\hat{prestige}$ = 20.8035385 + 0.0032691*6000 + 38.1200003*0 - 0.0024239(6000*0) + $\epsilon$
		= 20.8035385 + 19.6146
		= 40.4181385. \vspace{.5cm}
		
	The expected prestige score for a professional earning 6000 Canadian dollars (i.e., \textit{professional=1} and \textit{income = 6000}) is:
	
	$\hat{prestige}$ = 20.8035385 + 0.0032691*6000 + 38.1200003*1 - 0.0024239(6000*1) + $\epsilon$
	= 20.8035385 + 19.6146 + 38.1200003 - 14.5434
	= 63.9947388. \vspace{.5cm}
	
	63.9947388 - 40.4181385 = \textbf{23.5766003}
	
	On average, in comparison with a non-professional worker earning 6000 Canadian dollars, a professional worker earning 6000 Canadian dollars would be expected to have a prestige score that is 23.5766003 points higher. 
	
\end{enumerate}

\newpage

\section*{Question 2: Political Science}
\vspace{.25cm}
\noindent 	Researchers are interested in learning the effect of all of those yard signs on voting preferences.\footnote{Donald P. Green, Jonathan	S. Krasno, Alexander Coppock, Benjamin D. Farrer,	Brandon Lenoir, Joshua N. Zingher. 2016. ``The effects of lawn signs on vote outcomes: Results from four randomized field experiments.'' Electoral Studies 41: 143-150. } Working with a campaign in Fairfax County, Virginia, 131 precincts were randomly divided into a treatment and control group. In 30 precincts, signs were posted around the precinct that read, ``For Sale: Terry McAuliffe. Don't Sellout Virgina on November 5.'' \\

Below is the result of a regression with two variables and a constant.  The dependent variable is the proportion of the vote that went to McAuliff's opponent Ken Cuccinelli. The first variable indicates whether a precinct was randomly assigned to have the sign against McAuliffe posted. The second variable indicates
a precinct that was adjacent to a precinct in the treatment group (since people in those precincts might be exposed to the signs).  \\

\vspace{.5cm}
\begin{table}[!htbp]
	\centering 
	\textbf{Impact of lawn signs on vote share}\\
	\begin{tabular}{@{\extracolsep{5pt}}lccc} 
		\\[-1.8ex] 
		\hline \\[-1.8ex]
		Precinct assigned lawn signs  (n=30)  & 0.042\\
		& (0.016) \\
		Precinct adjacent to lawn signs (n=76) & 0.042 \\
		&  (0.013) \\
		Constant  & 0.302\\
		& (0.011)
		\\
		\hline \\
	\end{tabular}\\
	\footnotesize{\textit{Notes:} $R^2$=0.094, N=131}
\end{table}

\vspace{.5cm}
\begin{enumerate}
	\item [(a)] Use the results from a linear regression to determine whether having these yard signs in a precinct affects vote share (e.g., conduct a hypothesis test with $\alpha = .05$).
	
	\newpage		
	\item [(b)]  Use the results to determine whether being
	next to precincts with these yard signs affects vote
	share (e.g., conduct a hypothesis test with $\alpha = .05$).
	
	\vspace{7cm}
	\item [(c)] Interpret the coefficient for the constant term substantively.
	\vspace{7cm}
	
	\item [(d)] Evaluate the model fit for this regression.  What does this	tell us about the importance of yard signs versus other factors that are not modeled?
	
\end{enumerate}  


\end{document}
